\documentclass[12pt]{article}
\newcommand{\DocShort}{Gráfico comparativo}
\input{preamble_population.tex}
\begin{document}
\doccover{Método gráfico comparativo}{2026-05}
\handoutintro{Proyectar la población de una ciudad comparando su trayectoria con ciudades análogas ya desarrolladas.}{Cuando la ciudad bajo estudio tiene pocos datos o atraviesa una etapa de crecimiento similar a la de otras ciudades de referencia.}

\section{Datos ejemplo}
Ciudad objetivo X:
\begin{center}
\begin{tabular}{lcccc}
\toprule
Año & 1990 & 2000 & 2010 & 2020 \\
\midrule
Población & 28,000 & 34,000 & 41,000 & 49,000 \\
\bottomrule
\end{tabular}
\end{center}

Ciudades análogas:
\begin{center}
\begin{tabular}{lcccc}
\toprule
Ciudad & Punto 1 & Punto 2 & Punto 3 & Punto 4 \\
\midrule
A & 46,000 (1980) & 57,000 (1990) & 69,000 (2000) & 83,000 (2010) \\
B & 44,000 (1982) & 54,000 (1991) & 65,000 (2001) & 77,000 (2011) \\
C & 45,000 (1984) & 53,000 (1992) & 61,000 (2002) & 72,000 (2012) \\
D & 43,000 (1981) & 50,000 (1990) & 58,000 (2000) & 67,000 (2010) \\
\bottomrule
\end{tabular}
\end{center}

\section{Ecuaciones mínimas}
Tampoco es un método algebraico cerrado. La operación consiste en superponer trayectorias análogas y prolongar la ciudad objetivo siguiendo la curva media de referencia.

\section{Procedimiento}
\begin{enumerate}
\item Graficar las ciudades análogas.
\item Trazar una tendencia media visual.
\item Graficar la ciudad objetivo X en la misma escala.
\item Extender la curva de X imitando la forma de la trayectoria media.
\item Leer los valores futuros.
\end{enumerate}

\section{Ejemplo resuelto}
\begin{center}
\begin{tikzpicture}
\begin{axis}[
width=0.78\textwidth,
height=7.2cm,
xlabel={Año},
ylabel={Población},
xmin=1980, xmax=2040,
ymin=20000, ymax=90000,
grid=both,
major grid style={gray!25},
minor grid style={gray!10},
xtick={1980,1990,2000,2010,2020,2030,2040},
ytick={20000,30000,40000,50000,60000,70000,80000,90000},
scaled y ticks=false,
xticklabel style={/pgf/number format/1000 sep={}},
legend style={at={(1.02,1)}, anchor=north west, fill=white}
]
\addplot[color=Gray, mark=*] coordinates {(1980,46000) (1990,57000) (2000,69000) (2010,83000)};
\addlegendentry{A}
\addplot[color=Brown, mark=square*] coordinates {(1982,44000) (1991,54000) (2001,65000) (2011,77000)};
\addlegendentry{B}
\addplot[color=OliveGreen, mark=triangle*] coordinates {(1984,45000) (1992,53000) (2002,61000) (2012,72000)};
\addlegendentry{C}
\addplot[color=Orange, mark=diamond*] coordinates {(1981,43000) (1990,50000) (2000,58000) (2010,67000)};
\addlegendentry{D}
\addplot[color=medblue, mark=*, very thick] coordinates {(1990,28000) (2000,34000) (2010,41000) (2020,49000) (2030,59000) (2040,70000)};
\addlegendentry{X}
\addplot[densely dashed, color=teal, thick] coordinates {(1980,44500) (1990,53500) (2000,63250) (2010,74750) (2020,85000) (2030,95500)};
\addlegendentry{Tendencia media}
\node[anchor=south east, fill=white] at (axis cs:2028.6,61500) {59,000};
\node[anchor=south east, fill=white] at (axis cs:2038.6,72500) {70,000};
\end{axis}
\end{tikzpicture}
\end{center}

La lectura guiada del gráfico comparativo da:
\[
P_{2030}\approx 59,000,\qquad P_{2040}\approx 70,000
\]

\resultbox{La ciudad X se proyecta en 59,000 habitantes para 2030 y 70,000 para 2040.}

\section{Teoría breve}
La lógica del método es analógica: se asume que ciudades con condiciones semejantes tienden a recorrer trayectorias de crecimiento comparables.

\section{Observaciones de uso}
\begin{itemize}
\item La selección de ciudades comparables es la decisión crítica.
\item Las ciudades o comunidades análogas elegidas deben tener población mayor que la ciudad objetivo.
\item Deben parecerse en función económica, accesibilidad, jerarquía urbana y contexto regional.
\end{itemize}
\end{document}
